\chapter{Conclusiones y recomendaciones}
\label{ch:conclusiones}

% -----------------------------------------------------------------------
\section{Conclusiones}
\label{sec:conclusiones}

% TODO: una conclusión por objetivo específico (OE1..OE6), cada una
% respaldada por un resultado numérico del capítulo de resultados, y el
% veredicto sobre la hipótesis (error relativo medio obtenido vs 10 %).
% Evitar conclusiones genéricas: cada frase debe poder rastrearse a una
% tabla o figura.

% -----------------------------------------------------------------------
\section{Recomendaciones y trabajo futuro}
\label{sec:recomendaciones}

% TODO candidatos (depurar según lo que quede pendiente de verdad):
% - Integración Bender reproducible (fork de pulp_soc apuntando el core
%   modificado, hoy el clasificador se aplica sobre el checkout .bender).
% - Separación del consumo del core vs plataforma completa (medición en
%   el riel interno de la FPGA en lugar del jack de 5 V).
% - Extensión a otras configuraciones: otras frecuencias, otros cores
%   (CV32E40X/CVA6), traslado a ASIC.
% - Estimación en línea dentro del propio firmware (hoy el cálculo es
%   externo).

%%% Local Variables:
%%% mode: latex
%%% TeX-master: "main"
%%% End:
